% © 2026 Ing. David Jaroš — CC BY-NC-ND 4.0
%
% This work is licensed under a Creative Commons Attribution-NonCommercial-NoDerivatives
% 4.0 International License (CC BY-NC-ND 4.0).
%
% License History: Earlier drafts (up to v0.3) were released under CC BY 4.0.
% From v0.4 onward, all material is released under CC BY-NC-ND 4.0 to protect
% the integrity of the theoretical work during ongoing academic development.

\documentclass[11pt,a4paper]{article}
\usepackage{amsmath, amssymb, amsthm}
\usepackage{hyperref}
\usepackage{geometry}
\geometry{margin=2.5cm}

\newtheorem{theorem}{Theorem}
\newtheorem{proposition}{Proposition}
\newtheorem{definition}{Definition}
\newtheorem{remark}{Remark}
\newtheorem{conjecture}{Conjecture}

\title{Fermion Mass Generation in Unified Biquaternion Theory\\
\large Coupling Between \texorpdfstring{$\Theta$}{Theta} Modes, Mass Matrix, and Hierarchy}
\author{Ing.\ David Jaroš}
\date{2026}

\begin{document}
\maketitle

\begin{abstract}
We investigate the mechanisms generating particle masses within the
Unified Biquaternion Theory (UBT).  Starting from the Yukawa-type coupling
between KK modes of the $\Theta$-field, we derive the effective mass matrix
and identify the open gaps preventing a complete first-principles mass
derivation.  The comparison of the resulting structure with the Standard
Model Yukawa matrix is analysed.  All claims are labelled with their proof
status.
\end{abstract}

\tableofcontents

% ======================================================================
\section{Framework: Mass from the \texorpdfstring{$\Theta$}{Theta}-Field Action}
\label{sec:framework}
% ======================================================================

\subsection{Source of mass in UBT}

Particle masses in UBT originate from the potential term $V(\Theta)$ in the
canonical action (see \texttt{research/theta\_dynamics\_equations.tex} §4):
\begin{equation}
\label{eq:yukawa_lag}
  \mathcal{L}_{\mathrm{mass}}
  = -y_{ij}\,\bar\psi_{L,i}\,\phi\,\psi_{R,j} + \mathrm{h.c.}
    - \tfrac{1}{2} M^2_{ab}\,\Theta_a\Theta_b + \ldots,
\end{equation}
where:
\begin{itemize}
  \item $\phi$ is the scalar (Higgs) component of $\Theta$ (the $n=0$ KK zero-mode);
  \item $\psi_{L,i}$, $\psi_{R,j}$ are the left- and right-handed fermion
        components of the $n=i$ and $n=j$ winding modes (see
        \texttt{research/theta\_fermion\_emergence.tex} §4);
  \item $y_{ij}$ is the Yukawa coupling matrix;
  \item $M^2_{ab}$ is the scalar mass matrix from the Mexican-hat potential.
\end{itemize}

\subsection{Higgs mechanism and fermion mass}

After the vacuum condensation $\langle\phi\rangle = v$, the Yukawa coupling
produces:
\begin{equation}
\label{eq:mass_from_yukawa}
  m_{ij} = y_{ij}\,v.
\end{equation}

\begin{remark}[Status]
  Equation \eqref{eq:mass_from_yukawa} is the standard Yukawa mechanism
  applied to the UBT field structure.  Status: \textbf{[L1]}.
  The open problem is deriving $y_{ij}$ and $v$ from $R_\psi$ and the
  modular parameter $\tau_{\mathrm{mod}}$ without empirical fitting.
\end{remark}

% ======================================================================
\section{Mode Coupling and Mass Matrix}
\label{sec:coupling}
% ======================================================================

\subsection{KK mode overlap integrals}

The Yukawa coupling matrix element $y_{ij}$ arises from the overlap integral
of KK mode wave-functions on the $\psi$-circle:
\begin{equation}
\label{eq:overlap}
  y_{ij}
  = y_0 \int_0^{2\pi R_\psi} \!\! d\psi\;
    f_i^*(\psi)\, f_0(\psi)\, f_j(\psi),
\end{equation}
where $f_n(\psi) = (2\pi R_\psi)^{-1/2}\, e^{in\psi/R_\psi}$ are the
normalised KK wave-functions and $y_0$ is an overall coupling constant.

Using orthonormality:
\begin{equation}
\label{eq:y_selection}
  y_{ij} = y_0\, \delta_{i+j, 0 \bmod 1}.
\end{equation}

\begin{remark}[Status]
  Selection rule \eqref{eq:y_selection} is an exact algebraic identity
  (Status: \textbf{[L0]}).  It implies that KK modes $i$ and $j$ couple
  to the Higgs zero-mode only when $i = -j$ (or equivalently, when the
  pair is a winding/anti-winding mode pair).  This is consistent with
  charge conjugation symmetry.
\end{remark}

\subsection{Physical mass matrix}

For the three lightest winding modes $n = 1, 2, 3$:
\begin{equation}
\label{eq:mass_matrix}
  (M_f)_{ij} = y_{0}\,v\,\delta_{ij}
    + \text{(mixing from higher-order \textit{V} terms)}.
\end{equation}

At leading order (no mixing), the three generation masses are:
\begin{equation}
\label{eq:gen_masses}
  m_n = y_0\, v, \quad n = 1, 2, 3.
\end{equation}

\begin{remark}[Leading-order degeneracy]
  At leading order, \eqref{eq:gen_masses} predicts equal masses for all
  three generations — in contradiction to experiment.  The observed
  mass hierarchy $m_e : m_\mu : m_\tau \approx 1 : 207 : 3477$ must
  therefore arise from:
  \begin{enumerate}
    \item Off-diagonal mixing via higher-order $V$ terms, or
    \item Winding-number-dependent coupling constants $y_n$, or
    \item Holonomy profiles on the internal torus (see below).
  \end{enumerate}
  This is an identified open problem.  Status: \textbf{[O]}.
\end{remark}

% ======================================================================
\section{Holonomy Profile Mechanism}
\label{sec:holonomy}
% ======================================================================

\subsection{Higgs holonomy on the torus}

Following \texttt{unified\_biquaternion\_theory/fermion\_mass\_derivation\_complete.tex}
§Yukawa, the Higgs field on the internal torus $\mathbb{T}^2 = S^1_\chi \times S^1_\xi$
acquires a non-trivial holonomy profile:
\begin{equation}
\label{eq:holonomy}
  \Phi_{ij}(\chi, \xi) = \sum_{k \in \mathcal{H}} b_k\, e^{2\pi i k \cdot y},
\end{equation}
where $\mathcal{H} \subset \mathbb{Z}^2$ is the support set and
$b_k \in \{1, -1, i, -i\}$.

The effective Yukawa coupling for the $n$-th generation is then:
\begin{equation}
\label{eq:yukawa_holonomy}
  y_n = y_0 \int_{T^2} \Phi_{nn}(\chi,\xi)\, d\chi\, d\xi.
\end{equation}

\begin{remark}[Status]
  The holonomy mechanism is a motivated conjecture (\textbf{[L2]}).
  The discrete structure of $\mathcal{H}$ is motivated by gauge covariance
  (Status: \textbf{[L1]}), but its relation to the observed mass ratios
  is \textbf{semi-empirical [SE]} — the set $\mathcal{H}$ is fitted to
  reproduce the lepton masses.
\end{remark}

% ======================================================================
\section{Comparison with Standard Model Structure}
\label{sec:sm}
% ======================================================================

\subsection{Standard Model Yukawa matrix}

The SM Yukawa matrix in the mass eigenstate basis is diagonal:
\begin{equation}
  Y_{\mathrm{SM}} = \mathrm{diag}(y_e, y_\mu, y_\tau)
  \approx \mathrm{diag}(2.9\times10^{-6},\; 6.1\times10^{-4},\; 1.0\times10^{-2}).
\end{equation}

\subsection{UBT prediction and gaps}

\begin{center}
\begin{tabular}{lll}
\hline
Quantity & UBT status & Source \\
\hline
Higgs mass $m_h = \sqrt{2\lambda}\,v$ & \textbf{[SE]} — $\lambda$, $v$ fitted &
  \texttt{canonical/interactions/higgs.tex} \\
Electron mass $m_e \approx 0.511$~MeV & \textbf{[SE]} — 0.22\% accuracy &
  \texttt{consolidation\_project/electron\_mass/} \\
Muon–electron ratio $m_\mu/m_e \approx 207$ & \textbf{[O]} — not derived &
  \texttt{research/generation\_structure.tex} \\
Tau–muon ratio $m_\tau/m_\mu \approx 17$ & \textbf{[O]} — not derived &
  \texttt{research/generation\_structure.tex} \\
CKM matrix & \textbf{[P]} — partial & \texttt{appendix\_QA2\_quarks\_CKM.tex} \\
Neutrino masses & \textbf{[P]} — not treated & \texttt{STATUS\_NEUTRINOS.md} \\
\hline
\end{tabular}
\end{center}

% ======================================================================
\section{Towards a First-Principles Derivation}
\label{sec:first_principles}
% ======================================================================

\subsection{B\_m from action parameters}

The fermion mass formula (from \texttt{research/fermion\_mass\_program.md}) is:
\begin{equation}
  m(n) = A\, n^p - B_m\, n\,\ln(n),
\end{equation}
where $B_m \approx -14.099$~MeV is fitted.  A first-principles derivation
of $B_m$ requires identifying it as a combination of $R_\psi$, $y_0$, $v$,
and $\lambda$ from the action.

\begin{remark}[Status: Open]
  No complete derivation of $B_m$ from action parameters exists.
  Status: \textbf{[O]}.  See \texttt{research/fermion\_mass\_program.md}
  §Open Problem 1.
\end{remark}

\subsection{Strategy: moduli stabilization}

A natural path to fixing all mass parameters is moduli stabilisation:
if $R_\psi$ is fixed by a dynamical mechanism (e.g., from the modular
potential discussed in \texttt{research/modular\_dynamics.tex}), then
$v \propto R_\psi^{-1}$ and the entire mass spectrum could follow.

This represents the most important single open problem in the UBT mass
programme.

% ======================================================================
\section{Cross-References}
\label{sec:xref}
% ======================================================================

\begin{tabular}{ll}
\hline
Topic & File \\
\hline
Yukawa structure & \texttt{consolidation\_project/masses/yukawa\_structure.tex} \\
Fermion mass derivation & \texttt{unified\_biquaternion\_theory/fermion\_mass\_derivation\_complete.tex} \\
Electron mass & \texttt{consolidation\_project/electron\_mass/} \\
Three generations (proof) & \texttt{DERIVATION\_INDEX.md} \\
Generation mass ratios & \texttt{research/generation\_structure.tex} \\
Modular dynamics & \texttt{research/modular\_dynamics.tex} \\
\hline
\end{tabular}

\end{document}
